\chapter{矫正效果评估}
\label{chap:performance}
\newcommand{\qualitypending}{未测量}

本章采用标准斜边与网格图卡，评价矫正后的清晰度、几何直线度及边缘完整性，
对应赛题关于 MTF50 或 SFR、低直线度误差和无明显拉伸、断裂的要求。
本章规定同源比较、测试分组、指标计算和结果判定方法。
实物图卡测量数据尚未形成，结果表统一标记为“未测量”，不据此作出画质达标结论。
实时帧率、处理延迟与资源效率见第~\ref{chap:verification}~章。

\section{评价目标与测试原则}
\subsection{评价对象及指标}
矫正效果从三个方面评价：
\begin{enumerate}
\item 清晰度：从斜边图卡提取空间频率响应 SFR，以响应降至低频归一化值
50\% 时的频率 MTF50 作为主要数值，同时保留完整曲线和中心、边缘差异。
\item 直线度：从网格图卡提取长直线，统计矫正前后相对于最佳拟合直线的
RMS、最大和 P95 正交距离，重点检查画面外围的残余弯曲。
\item 边缘完整性：检查有效区域内图形的局部尺度、连通性和跨块连续性，
记录异常拉伸、断线、错位、锯齿及其位置，辅以局部放大图。
\end{enumerate}

赛题未给出统一的 MTF50、直线度或拉伸误差数值门限，
评价门限由图卡、分辨率、有效视场和测量不确定度共同确定，
采用预先固定的内部判据，与官方要求分别列示。
“直线度误差极低”需要同时给出像素单位残差和相对于测试线长度的比例，
“边缘无明显异常”则需要有明确观察倍率、有效区域和辅助量化记录。

\subsection{同源比较与控制变量}
主测试固定曝光、增益、焦距、对焦、相机位置、图卡距离、照明和标定参数。
优先保存同一源帧的输入图像与 FPGA 输出；若只能分别采集，
应使用静止图卡和固定姿态，并记录配对方式及可能的采集差异。
输出使用原生 $1920\times1080$ 图像，保存无损文件，
关闭采集软件中的缩放、锐化、降噪和有损压缩。
显示器截图、拍屏图和报告排版后的缩略图不作为主要定量分析输入。

默认实时路径采用 RGB565 存储及最近邻采样。
为分离颜色量化与几何重采样的影响，可另保存按相同规则量化、再展开的输入对照，
并建立同标定、同姿态、同输出视场的独立浮点参考。
几何主测试关闭自适应 Gamma；增强功能开启的测试单独分组，
避免将对比度变化或锐化造成的响应变化解释为几何矫正带来的清晰度提升。
\sourcefile{\path{problem.md} 中“矫正效果评估”要求；
\path{../Pgl50h/src/video/pose/pose_renderer_tiled.v}；
\path{../Pgl50h/src/video/format/rgb888_to_rgb565.v}}

\section{图卡、采集环境与测试分组}
\subsection{标准测试图卡与布置}
清晰度测试采用斜边图卡或符合相应测量规范的倾斜黑色方块，
在近水平、近垂直方向各提供一条略带倾角的高对比边缘，
建议边缘相对像素行或列偏转约 $5^\circ$，以取得足够的亚像素采样相位。
图卡规格、测量规范和分析工具共同构成测试条件；
斜边倾角只用于取得有效采样，不单独作为符合测量规范的判据。
边缘应清晰、无反光和明显印刷毛刺，亮暗两侧均避免饱和。

几何测试采用已知线宽和间距的平面网格图卡，覆盖画面中心、四边和四角。
图卡应固定在平整基底上，记录制造精度与平整度，
以免将打印误差、纸张弯曲或安装倾斜计入算法残差。
倾斜方块同时用于观察边缘连续性，但清晰度所用的轻微边缘倾角，
与系统 yaw/pitch 姿态角是不同条件，应分别记录。

\begin{table}[htbp]
\centering\small
\caption{图卡及采集条件记录}\label{tab:quality-setup}
\begin{tabularx}{\textwidth}{>{\raggedright\arraybackslash}p{3.5cm}Y}
\toprule
记录项目 & 实际配置 \\
\midrule
相机、镜头与采集设备 & USB 工业相机与 HDMI 采集卡；测量配置未记录 \\
图卡与测量规范 & 斜边与网格图卡；规格及规范未记录 \\
距离、焦距及对焦 & 未记录 \\
曝光、增益与照明 & 未记录 \\
标定与安装关系 & 采用系统标定模型；测量安装条件未记录 \\
FPGA 实现与处理配置 & 最近邻分块采样；测量配置文件未记录 \\
图像与有效区域 & $1920\times1080$；实际测量格式未记录 \\
分析环境与重复采样 & 未记录 \\
\bottomrule
\end{tabularx}
\end{table}

\subsection{模式与姿态测试矩阵}
表~\ref{tab:quality-modes} 的模式用于区分各项处理的影响。
主结果应同时注明三开关掩码、姿态和 Gamma 状态，不能仅写“开启矫正”。

\begin{table}[htbp]
\centering\small
\caption{矫正效果的测试分组}\label{tab:quality-modes}
\begin{tabularx}{\textwidth}{>{\raggedright\arraybackslash}p{1cm}>{\raggedright\arraybackslash}p{4cm}Y}
\toprule
编号 & 处理条件 & 评价目的 \\
\midrule
A & 原始输入与同格式量化对照 & 保留未经几何处理的成像基准，识别 RGB565 量化影响。\\
B & 镜头矫正开启，姿态/Gamma 关闭 & 测量镜头去畸变后的直线度与清晰度变化。\\
C & 镜头、姿态矫正开启，Gamma 关闭 & 作为联合矫正主测试，评价不同姿态下的几何和边缘质量。\\
D & 镜头、姿态及自适应 Gamma 开启 & 补充正常、暗光、强光场景，单独评价增强后的观感和响应。\\
\bottomrule
\end{tabularx}
\end{table}

静态姿态建议覆盖零姿态、yaw/pitch 各自正负边界、双轴边界组合及代表性中间角度，
边界范围以实际采用配置为准，默认约为各轴 $\pm15^\circ$。
机械角度、IMU 读数与渲染器实际锁存角度应同时记录。
每个状态建议保存不少于 10 个完整输出帧，统计逐帧结果与重复性；
该采样数量用于规定测量方法，不计入已完成的测试覆盖。
另设缓慢连续转动和功能切换用例，检查运动中是否出现局部撕裂、跳线或边缘断续，
其结果与静止图卡测量分开记录。

\subsection{中心、边缘与角部 ROI}
将有效图像划分为中心、上/下/左/右边缘和四角共九个位置，
每处选择合适的斜边 ROI；网格则保留跨越较大范围的完整长线。
ROI 编号、坐标、大小、边缘方向和对应物理图卡位置应固定保存。
矫正前后同一物体可能位于不同像素位置，应按图卡特征或独立几何映射匹配，
不能用相同像素坐标框代替同源区域匹配。

\begin{figure}[htbp]
\centering
\begin{tikzpicture}[x=1cm,y=1cm,font=\small]
\draw[thick,TemplateBlue] (0,0) rectangle (10.8,5.2);
\foreach \x/\y/\name in {1.3/4.3/TL,5.4/4.3/T,9.5/4.3/TR,
  1.3/2.6/L,5.4/2.6/C,9.5/2.6/R,1.3/0.9/BL,5.4/0.9/B,9.5/0.9/BR} {
  \draw[TemplateBlue,fill=TemplateBlue!7] (\x-0.65,\y-0.35) rectangle (\x+0.65,\y+0.35);
  \node at (\x,\y) {\name};
}
\end{tikzpicture}
\caption{九位置 ROI 布置示意。位置框只表示选点方式，实际 ROI 需落在可测且有效的图卡边缘上。}
\label{fig:quality-roi}
\end{figure}

对于姿态矫正后移出有效视场的位置，记录为“不适用”并说明原因，
不能填为零或在汇总时悄然删除。
黑边和越界填充依据几何有效掩码确定，不能仅按 RGB 为零判定，
因为图卡本身也包含黑色区域。
测试记录以源帧编号关联姿态、有效掩码与 ROI 清单。


\section{基于斜边的 SFR 与 MTF50 清晰度测量}
\subsection{测量方法与计算流程}
SFR 描述图像对不同空间频率的响应。由斜边估计得到的响应包含镜头、传感器、
相机处理、量化与重采样的共同影响，本章将其作为采集至 FPGA 输出链路的清晰度指标。
在未完成线性化和方法核验前，不将其解释为镜头独立的光学 MTF。

对每个 ROI 按以下流程处理，并保存中间数据：
\begin{enumerate}
\item 检查图像无裁剪饱和、明显反光和运动模糊，将像素转为固定定义的亮度或灰度。
若已知相机及输出响应，按实测 OECF 或适当的逆响应完成线性化；
若未知，则记录处理假设，将结果标为当前编码条件下的 SFR，不能直接套用未经核验的 Gamma。
\item 定位斜边并拟合其位置，将 ROI 像素投影到边缘法向，
按亚像素距离分箱形成边缘扩展函数 ESF。记录拟合方式、ROI 大小、倾角和过采样倍数。
\item 对 ESF 求导得到线扩展函数 LSF，按所选算法执行窗函数、
离散微分响应修正和频域归一化，记录所用实现及参数。
\item 计算 SFR 曲线，确定低频归一化响应首次向下穿过 0.5 的频率，
记录 MTF50，并保留边缘过冲、振铃及非单调响应的情况。
\item 对同组重复帧和各个位置重复计算，报告样本数、均值、标准差及最差位置。
剔除异常图像必须使用预先规定的规则，并保留原始帧和剔除原因。
\end{enumerate}

连续形式下的关系可写为
\begin{equation}
\operatorname{LSF}(x)=\frac{\mathrm{d}\operatorname{ESF}(x)}{\mathrm{d}x},\qquad
\operatorname{SFR}(f)=
\frac{\left|\mathcal{F}\{\operatorname{LSF}(x)\}(f)\right|}
{\left|\mathcal{F}\{\operatorname{LSF}(x)\}(0)\right|},
\end{equation}
\begin{equation}
f_{50}=\min\{f>0:\operatorname{SFR}(f)=0.5
\text{，且为首次下降交点}\}.
\label{eq:quality-mtf50}
\end{equation}
实际离散算法还需完成上述窗函数和微分修正，不能只对未经处理的阶跃图像直接做 FFT。
主报告在输出采样的 $0<f\le0.5$ cycles/pixel 范围内给出结果；
若该范围内没有有效 50\% 交点，应记录“范围内未交叉”或“无法可靠估计”，
不能外推填值。ESF 过采样改善估计的相位采样，并不等于输出图像具有更高的独立空间分辨率。

\subsection{频率单位、方向及局部尺度}
主表统一使用 cycles/pixel。若另以图像高度归一化为 LW/PH，须说明线宽口径，
按输出高度 $H$ 使用
\begin{equation}
\operatorname{MTF50}_{\rm LW/PH}=2Hf_{50},\qquad H=1080.
\end{equation}
H 方向的频率响应由近垂直斜边测得，V 方向由近水平斜边测得；
这里 H/V 指频率变化方向，而非边缘延伸方向。

几何矫正会改变局部放大率，直接比较矫正前后的 cycles/pixel 可能把尺度变化
误判为清晰度损失或提升。应记录对应图卡区域在所测方向上的 pixel/mm，
必要时换算为同一物方尺度的 cycles/mm 后比较，或按预先确定的方法统一局部尺度。
用于物方换算的尺度必须对应同一方向及局部区域。
任何额外缩放本身也会引入插值，应另存原生图测量结果并记录缩放方法。
Gamma 开启组单独记录亮度响应和当前曲线，避免与关闭组直接混合平均。

\subsection{清晰度结果记录}
每个姿态、模式和照明条件分别统计表~\ref{tab:quality-mtf} 的指标。
H/V 表示两个测量方向，数值与相应 ROI、曲线和局部尺度关联。
记录项包括测试组、源帧、处理模式、姿态、ROI 尺寸、样本数与线性化方法。

\begin{table}[htbp]
\centering\small
\caption{分区域 MTF50 结果记录（主单位：cycles/pixel）}\label{tab:quality-mtf}
\begin{tabularx}{\textwidth}{>{\raggedright\arraybackslash}p{1.4cm}YYYY}
\toprule
区域 & 原图 H/V & 矫正后 H/V & 局部尺度及比较口径 & ROI/曲线证据 \\
\midrule
中心 C & \qualitypending & \qualitypending & \qualitypending & \qualitypending \\
上 T & \qualitypending & \qualitypending & \qualitypending & \qualitypending \\
下 B & \qualitypending & \qualitypending & \qualitypending & \qualitypending \\
左 L & \qualitypending & \qualitypending & \qualitypending & \qualitypending \\
右 R & \qualitypending & \qualitypending & \qualitypending & \qualitypending \\
左上 TL & \qualitypending & \qualitypending & \qualitypending & \qualitypending \\
右上 TR & \qualitypending & \qualitypending & \qualitypending & \qualitypending \\
左下 BL & \qualitypending & \qualitypending & \qualitypending & \qualitypending \\
右下 BR & \qualitypending & \qualitypending & \qualitypending & \qualitypending \\
\bottomrule
\end{tabularx}
\end{table}

清晰度统计包含逐帧数据、各 ROI 的均值与标准差、外围最差值及测量不确定度。
分析应结合完整 SFR、ESF 和局部边缘图，说明变换尺度、最近邻锯齿及量化的影响，
避免仅用单个 MTF50 数值断言视觉质量全面提高。


\section{网格直线度与残余畸变}
\subsection{直线提取及误差定义}
对网格中的长直线进行亚像素边缘或中心线提取。
宽黑线宜使用两侧边缘确定中心线，并固定算法和阈值，
减少线宽及明暗阈值变化对位置估计的影响。
每条测试线应尽量覆盖较长有效范围，在中心、四边和角部选择横、纵及必要的斜向线。
ROI、端点排除距离、拟合点数、异常点处理和遮挡排除规则应提前确定。

采用正交最小二乘等适合任意方向的拟合方式，得到
\begin{equation}
ax+by+c=0,\qquad a^2+b^2=1.
\end{equation}
对该线上 $n$ 个提取点 $(x_i,y_i)$，定义带符号残差及距离
\begin{equation}
r_i=ax_i+by_i+c,\qquad d_i=|r_i|,
\end{equation}
并计算
\begin{equation}
D_{\rm RMS}=\sqrt{\frac{1}{n}\sum_{i=1}^{n}r_i^2},\qquad
D_{\max}=\max_i d_i,\qquad
D_{95}=\operatorname{percentile}_{95}(d_i).
\end{equation}
若有效测试线段长度为 $L$ 个输出像素，归一化最大直线度误差为
\begin{equation}
\varepsilon_{\rm line}=\frac{D_{\max}}{L}\times100\%.
\label{eq:quality-linearity}
\end{equation}
报告中同时给出像素单位和归一化值，并保留沿线位置对应的带符号残差曲线，
以识别连续弯曲和局部跳变。

透视投影下直线仍为直线，但平行线不一定在图像中保持平行，
线段倾斜也不等于残余畸变。
直线度评价不强行将所有线拟合为水平或垂直；
图卡姿态、校正后视场和投影约定造成的整体倾斜，应与残余弯曲分开说明。
图卡平整度、采集噪声及亚像素提取误差也应计入测量不确定度，
不能仅依据软件输出的小数位数宣称达到对应精度。

\subsection{直线度记录与判定}
每条线均保留原图和矫正后结果，按相同物理线段匹配。
当两者覆盖长度不同或局部尺度变化较大时，应说明像素残差的可比性，
必要时增加同尺度结果；不得通过缩短线段或删除外围采样点来降低残差。

\begin{table}[htbp]
\centering\small
\caption{直线度测量结果（距离单位：pixel）}\label{tab:quality-linearity}
\begin{tabularx}{\textwidth}{>{\raggedright\arraybackslash}p{2cm}YYYYY}
\toprule
线段位置 & 长度/点数 & 原图 RMS & 矫正 RMS & 矫正最大/P95 & 归一化最大误差 \\
\midrule
中央横线 & \qualitypending & \qualitypending & \qualitypending & \qualitypending & \qualitypending \\
中央纵线 & \qualitypending & \qualitypending & \qualitypending & \qualitypending & \qualitypending \\
上/下边缘线 & \qualitypending & \qualitypending & \qualitypending & \qualitypending & \qualitypending \\
左/右边缘线 & \qualitypending & \qualitypending & \qualitypending & \qualitypending & \qualitypending \\
角部代表线 & \qualitypending & \qualitypending & \qualitypending & \qualitypending & \qualitypending \\
\bottomrule
\end{tabularx}
\end{table}

上/下、左/右及四角实测时分别展开记录，不以一个平均值掩盖最差位置。
逐线统计包含 RMS、最大误差、P95、线段长度和最差位置。
直线度结论由测量值、不确定度与预设判据共同决定。


\section{边缘拉伸、断裂与局部完整性}
\subsection{异常拉伸与局部尺度}
使用网格相邻交点间距、方块边长和长宽比检查局部形变，
重点观察四角、姿态边界及有效视场外围。
几何矫正本身会产生预期的尺度变化，所以不能将所有拉伸都认定为缺陷，
也不能要求透视图像中所有网格间距天然相等。
应与相同视场、姿态和采样约定的独立参考，或具有明确成像几何的物方基准比较。

对匹配的局部线段，设输出长度为 $\ell$，参考长度为 $\ell^*$，定义
\begin{equation}
\varepsilon_{\ell}=\frac{|\ell-\ell^*|}{\ell^*}\times100\%.
\end{equation}
横向、纵向分别统计，必要时增加相对于参考长宽比的偏差。
参考线段、交点定位方法、尺度标定和允许误差均应记录。
若未建立独立几何参考，则只保留尺寸分布和观察结论，
不将其填作已经测得的“拉伸误差”。

\subsection{断裂、错位和跨块连续性}
在有效区域内，沿同一条图卡线或方块边缘跟踪连通性，
记录异常缺口数、最大缺口长度和跳变位置。
检测前应确定最小缺口长度、允许的像素连接方式及噪声处理参数，
并对算法检出的异常逐一查看原生像素，避免阈值噪声造成误报。
设计产生的越界黑边、图卡真实间断及被遮挡区域单独标注，不能计作处理断线。

默认输出块大小为 $128\times32$，应重点检查
$x=128k$、$y=32m$ 附近穿过网格线或斜边的位置，
以及图像末行、末列和最后一个不完整块。
对边界两侧比较线条位置、连续性及亮度，检查是否出现整行错位、缓存旧数据或块缝。
最近邻采样可能产生台阶与锯齿，应记录其尺度及可见程度，
并与缺失像素造成的真正断裂区分。

观察统一采用原生 1:1 显示，辅以注明倍率的整数倍最近邻放大，
避免展示软件平滑插值掩盖块缝和缺口。
动态检查则保存连续输出帧或视频，记录姿态、功能切换时刻和异常帧号。

\begin{table}[htbp]
\centering\small
\caption{边缘完整性测量项目}\label{tab:quality-edge}
\begin{tabularx}{\textwidth}{>{\raggedright\arraybackslash}p{3cm}YYY}
\toprule
检查项目 & 区域/条件 & 结果记录 & 局部图/帧号 \\
\midrule
局部横/纵尺度误差 & \qualitypending & \qualitypending & \qualitypending \\
方块长宽比偏差 & \qualitypending & \qualitypending & \qualitypending \\
异常断线数/最大缺口 & \qualitypending & \qualitypending & \qualitypending \\
跨块边界位置跳变 & \qualitypending & \qualitypending & \qualitypending \\
斜边台阶及可见锯齿 & \qualitypending & \qualitypending & \qualitypending \\
末行/末列/不完整块 & \qualitypending & \qualitypending & \qualitypending \\
姿态变化及开关切换 & \qualitypending & \qualitypending & \qualitypending \\
有效覆盖与越界区域 & \qualitypending & \qualitypending & \qualitypending \\
\bottomrule
\end{tabularx}
\end{table}


\section{综合判据与评价结论}
\subsection{门限与验收表}
表~\ref{tab:quality-acceptance} 中的门限应在正式测量前确定。
不同姿态与方向分别判断，主结论使用有效 ROI 中最差位置，
重复测量时同时考虑统计波动及测量不确定度。
若部分位置无法测量，应列出原因及有效视场覆盖，不能默认为通过。

\begin{table}[htbp]
\centering\small
\caption{矫正效果综合评价}\label{tab:quality-acceptance}
\begin{tabularx}{\textwidth}{>{\raggedright\arraybackslash}p{3.2cm}YYY}
\toprule
指标 & 预设判据 & 最差实测结果 & 结论/证据 \\
\midrule
MTF50 与完整 SFR & 未设定 & \qualitypending & \qualitypending \\
直线度 RMS & 未设定 & \qualitypending & \qualitypending \\
最大/归一化直线度 & 未设定 & \qualitypending & \qualitypending \\
局部尺度/长宽比偏差 & 未设定 & \qualitypending & \qualitypending \\
异常断裂与位置跳变 & 未设定 & \qualitypending & \qualitypending \\
原生倍率下边缘观感 & 未设定 & \qualitypending & \qualitypending \\
有效视场覆盖 & 未设定 & \qualitypending & \qualitypending \\
\bottomrule
\end{tabularx}
\end{table}

\subsection{评价结论与数据记录}
本章建立了斜边清晰度、网格直线度和边缘完整性三类评价方法，
分别对应频域响应、几何残差和局部连续性。
中心与外围、水平与垂直、静态与动态条件分组统计，
以有效区域内的最差位置和测量不确定度支撑综合判断。

实物图卡的同源输入输出、逐 ROI 曲线、直线残差及局部放大图尚未形成完整测量记录，
因此表~\ref{tab:quality-mtf}、表~\ref{tab:quality-linearity} 和
表~\ref{tab:quality-edge} 的结果列保留为“未测量”，
表~\ref{tab:quality-acceptance} 不作通过判定。
软件合成图卡与 RTL 功能测试分别验证分析方法和数据通路，
不替代真实相机与 FPGA 输出的定量画质测量。

每组测量采用源帧身份关联输入输出、量化对照、浮点参考、有效掩码和 ROI 坐标，
同时记录曝光、照明、姿态、功能配置、分析方法、原始数据与文件校验值，
保证结果可复核。
